\cotefile{44}
\scanpage{2}\mspage{\unsure{1|7}}
\begin{center}\uline{Programme de Travail}\end{center}

\deleted{\unsure{B|3}}\uline{A}

% Right-hand square bracket drawn to the left of Th. 1.
\noindent\llap{\smash{\raisebox{-2.2ex}{$\left.\rule{0pt}{4.8ex}\right]$}}\quad}\uline{Th. 1}\qquad $N(\unsure{\zeta|\xi})\subset\operatorname{Alb}(S)$ n'est autre \inserted{\deleted{des}ensemblistement} que le noyau de l'application naturelle
\[
\varphi_{\zeta}:\quad\operatorname{Alb}(\unsure{s|S})\longrightarrow
\mathcal{Z}^{i}(X)/\mathcal{Z}^{i}_{\mathrm{alb}}(X)
\]

% Two overlapping right brackets in the left margin.
\noindent\llap{\smash{\raisebox{-1.5ex}{$\left.\rule{0pt}{3.6ex}\right]$}}\quad}\llap{\smash{\raisebox{-3.3ex}{$\left.\rule{0pt}{3.1ex}\right]$}}\quad}Th. 2\qquad La dimension de $\operatorname{Alb}(S)/N(\zeta)=P_{\zeta}$ est majorée par un nb ne dépendant que de $X$ [en fait, par $\frac{1}{2}b^{2i-1}(X)$ !]

Ces deux th. impliquent

\noindent\llap{\smash{\raisebox{-2.2ex}{$\left.\rule{0pt}{4.8ex}\right]$}}\quad}\uline{Th. 3}\quad $Z^{i}_{\mathrm{alg}}(X)/\mathcal{Z}^{i}_{\tikz[baseline=(a.base)]{\node[ellipse,draw,inner sep=0.3pt] (a) {$\scriptstyle\mathrm{alb}$};}}(X)=P^{i}(X)$ a une structure naturelle de V.A., $\deleted{\unsure{P}}$et $P^{i}(X)$ et $P^{\unsure{2m-i|n+1-i|n-i}}(X)$ sont en dualité canonique

% slanted; circled note with an arrow to Th. 4
\margin{\tikz[baseline=(n.base)]{\node[draw,ellipse,inner sep=1pt,align=center] (n) {\unsure{Canulé|Annulé|Cancelé}\\tel quel};\draw[->] (n.south) -- ++(0.18,-0.2);}}
% slanted
\margin{\textquotedblleft{}\uline{Hodge}}
\uline{Th. 4}\quad Si $n=2m-1$, alors l'autodualité canonique de $P^m(X)$ [donc \unsure{la|sa} corr. sur $\underline{P}^m(X)$] est \unsure{associé|associée} à une \uline{polarisation} de $P^m(X)$
% archive sheet number: 1
\scanpage{3}
% A long left stroke and diagonals cancel the entire draft.
\deleted{Soit $A=\coprod_{i\geq0}A^i$ un anneau gradué, à degrés $\geq0$ et soit $\hat A=\prod_{i\geq0}A_i$ son complété pour la filtration canonique. Si cette dernière est \unsure{adique} et $A$ noth. [p.ex. $A_0$ noth., $A$ engendré sur $A_0$ par $A_1$, $A_1$ de type fini sur $A_0$] alors $\hat A$ est \deleted{\unsure{fid.}} plat sur $A$. Pour tout $A$-module $M$, nous posons
\[
\hat M=M\otimes_A\hat A
\]
Supposons que $M$ soit lui-même gradué, et de type fini, alors on a un isom naturel
\[
\hat M=\prod_i M_i
\]
Si $M$ satisfait seulement TF, alors il existe une suite exacte
\[
0\mathrel{\unsure{-|\to}}M''\mathrel{-}M\mathrel{-}M'\mathrel{-}0
\]
\deleted{\unsure{[$M'=$}} $M'$ à degrés majorés, $M''$ de type fini\inserted{($M''=\coprod_{i\geq\unsure{\nu|v}} M_{\unsure{\nu|i}}$)}, d'où
\[
0\mathrel{-}\hat M''\mathrel{-}\hat M\mathrel{-}\hat M'\mathrel{-}0
\]
Et on a
\[
\begin{aligned}
\hat M''&=\prod_{i\geq\nu} M''_i=\prod_{i\geq\unsure{0|\nu}} M_i\\
\hat M'&=M'
\end{aligned}
\]
comme on voit de suite [car $M'$ est \deleted{\illegible{1}} loc annulé par des idéaux de définition].

Donc $\unsure{\hat M^{\cdot}|\hat M}$ est formé de sommes $\sum_{i\in\unsure{Z|\mathbb{Z}}}x_i$ telles \deleted{\unsure{pour}} que $x_i\unsure{=|\ne}0$ pour $i\longrightarrow-\infty$. \deleted{\unsure{\illegible{1}|toujours cas}} \inserted{\uline{En tous cas}} on voit que $M\mathrel{\unsure{\leadsto|\mapsto}}\hat M$ sur les modules \uline{gradués} \deleted{\unsure{TF}} est exact et \uline{\unsure{fidèle|fidèl}}}
% archive sheet number: 2
\scanpage{4}\mspage{2}
\[
\deleted{2i-1+2j-1=2(i+j-1)=\unsure{2n|2m|n}}
\]

% slanted
\margin{\textquotedblleft{}\uline{Néron}\textquotedblright{}}
\uline{Th. 5}\qquad Soient $X$ variété projective n.s. de dim $n$ \deleted{\illegible{2}}, $Z^i,Z'^{j}\in\mathcal{Z}^{\deleted{\unsure{*|n}}}(X)$\inserted{ de codim $i,j$, avec \deleted{\unsure{\illegible{1}|codim}} $i+j=n$}. Supposons que $z=Z\mathbin{.}Y$ et $z'=Z'\mathbin{.}Y$ soient alg. équiv. à $0$\inserted{ sur $Y$} [$Y$ section hyperplane\inserted{ générique d'un pinceau\inserted{ dont \deleted{toutes} les fibres sont géom. intègres\ldots}}] ; ils sont de codim \deleted{\unsure{$m$ et $m$}}\inserted{$i,j$}, avec \deleted{$\dim Y=2m-1$}\inserted{$i+j=n=\dim Y+1$} donc \deleted{\unsure{ce sont des|définissent des}}\inserted{définissent} des éléments de $P^{\deleted{\unsure{m}}i}(Y),P^{j}(Y)$ et par la forme de Néron-Tate [définie grâce à l'accouplement entre $P^{i}(Y)$ et $P^{j}(Y)$] on \deleted{\unsure{a donc}} définit $\langle z,z'\rangle$. Ceci posé, on a \deleted{[au signe près ??]}
\[
\langle z,z'\rangle=(-1)^m Z\mathbin{.}Z'
\]

% slanted; circled
\margin{\uline{N.B} Sauf si $2i=2j=n$ les deux \unsure{mbs|mbres|membres} sont nuls par Lefschetz \unsure{aux|au|en} \unsure{deux genres|tous genres|sens grossier}}
% A triangular border encloses the note. An arrow points up to Ceci.
\margin{\tikz[baseline=(n.base)]{\node[text width=2cm,align=left,inner sep=1pt] (n) {Doit résulter d'un énoncé plus général sur des fibrations, et admet une \uline{variante} locale};\draw (n.south west) -- (n.north west) -- ([yshift=3mm]n.north east) -- (n.south east) -- cycle;\draw[->] (n.north) -- ++(0,0.3);}}

\deleted{\uline{\unsure{Cor.|Corr.}}} \uline{Th. 6}\qquad Sous les conditions du th 5,

% slanted
\margin{\textquotedblleft{}\uline{Lefschetz}\textquotedblright{}}
Si $2i\leq n-1$ \ $P^i(X)\mathrel{-}P^i(Y)$ est une \uline{isogénie} ; \deleted{\unsure{et}} \uline{si $2i=n$} le \uline{facteur fixe} de $P^i(Y)$ est identique à l'image de $P^i(X)$ ; \deleted{\unsure{[} \deleted{\unsure{et}} C'est égal à $P^i(X)$} \deleted{sauf [Si \deleted{\unsure{pour}} \unsure{$2i=n$} $n=2i$ \unsure{ou} $n$ pair et $i=\frac{n}{2}$]} \deleted{Le noyau de $P^i(X)\mathrel{-}P^i(Y)$ est fini si $2i\leq n-1$} [Si \deleted{$\unsure{2i|2}$} $2i\geq n+1$, $P^i(X)\mathrel{-}P^i(Y)$ est \inserted{surjectif}]
% archive sheet number: 3
\scanpage{5}
% One long oblique penstroke crosses the draft.
\deleted{Considérons les objets suivants

1) \quad $(P,\deleted{\unsure{\Pi|T}}\Delta)$, où $P$ est un groupe abélien libre de rang fini, $\deleted{\Pi}\Delta$ une partie finie de $P$, satisfaisant les conditions

a) \quad $\deleted{\Pi}\Delta$ engendre $P$

b) \quad Si $a\in\deleted{\Pi}\Delta$, on a $-a\in\deleted{\Pi}\Delta$, mais $na\notin\deleted{\Pi}\Delta$ si $n\ne1,-1$, $n\in\mathbb{Q}$

c) \quad Pour tout $a\in\deleted{\Pi}\Delta$, il existe une \textquotedblleft{}symétrie\textquotedblright{} $S_a$ appliquant $a$ sur $-a$, et telle que $S_a(\deleted{\Pi}\Delta)\subset\deleted{\Pi}\Delta$.

N.B. Alors les $S_a$ en question sont uniques, \deleted{\unsure{car en}}\inserted{étant en} effet \deleted{\unsure{induit}} l'unique symétrie euclidienne\inserted{ de $P^{\mathbb Q}$} pour la forme quad. suivante : $\deleted{\unsure{(x,y)=(x',y')=}}$ celle déduite par dualité de la forme $(x,y)=\sum_{a\in\Delta}\langle x,a\rangle\langle y,a\rangle$ sur $(P^{\mathbb Q})'=(P')^{\mathbb Q}$

On aura alors
\[
S_{\alpha}x=x-\langle\unsure{r|\gamma}_{\alpha},x\rangle\alpha
\]
pour un $r_{\alpha}\in P'^{\mathbb Q}$ bien déterminé. On montre que l'on a en fait $r_{\alpha}\in P'$, i.e. les $\deleted{\unsure{r}}$
\[
c_{\alpha\beta}=\unsure{-|{}}\langle r_{\alpha},\beta\rangle\qquad\text{sont }\in Z
\]
[proviendra du fait que chaque $\alpha$ est non divisible dans $P$, et ceci du fait que $\alpha$ fait partie d'un \textquotedblleft{}système fondamental de racines\textquotedblright{}]

2) \quad $(P,P',\Delta,\Delta',\varphi)$, où $P,P'$ sont deux groupes abéliens libres de rang fini \deleted{\unsure{en dualité}}\inserted{\unsure{accouplés à} $Z$}, $\Delta\subset P$, $\Delta'\subset P'$, $\varphi$ une bijection de $\Delta$ sur $\Delta'$, \deleted{\unsure{et}} les conditions \deleted{suivantes} étant vérifiées :

a) \quad $P^{\mathbb Q}$ et $P'^{\mathbb Q}$ sont en dualité, $\Delta$ engendre $P$, \inserted{$\Delta'$ engendre $P'$}

\margin{Si $\alpha'=\varphi(\alpha)$, on a $\langle\alpha',\alpha\rangle=2$}
b) \quad Pour tout $\deleted{\unsure{-}}\alpha'\in\Delta'$, \inserted{$\alpha\in\Delta$ tels que $\varphi(\alpha)=\alpha'$,} l'endomorphisme
\[
S_{\alpha}:\quad x\longmapsto x-\langle\alpha',x\rangle\alpha
\]
applique $\Delta$ dans $\Delta$, \deleted{\illegible{1}}

\noindent$\Updownarrow$\qquad b$'$) Pour tout $r\in\Delta'$, $\alpha\in\Delta$ tels que $\varphi(\alpha)=r$,
\[
S'_{\unsure{\alpha'|\alpha}}:\quad x'\longmapsto x'-\langle x',\alpha\rangle\alpha',\quad\text{invarie }\Delta'
\]}
% archive sheet number: 4
\scanpage{6}\mspage{3}
% slanted
\margin{\textquotedblleft{}\uline{\unsure{Matsusaka}}\textquotedblright{}}
\uline{Th. 7}\qquad L'équivalence numérique $=$ $\tau$-équivalence \deleted{(}[cycles algébriques sur $X$ projective non singulière\deleted{)}].
% archive sheet number: 5
\scanpage{7}\mspage{- d -}
\begin{center}\uline{Commentaires}\end{center}

Le rédacteur a suivi aussi fidèlement que possible les instructions de La Tribu n\textsuperscript{o} 57. Le plan n'a été modifié qu'en ce qui concerne les remords-fourre-tout du début, que le rédacteur a jugé préférable de grouper en 3 paragraphes préliminaires au lieu d'un, le 3e étant consacré aux lemmes sur les graphes qui traînaient un peu partout, et dont en fait on a besoin dès le début, (p. ex. pour le critère pour qu'une matrice de Coxeter soit connexe). Conformément aux injonctions de La Tribu, le rédacteur a suivi la terminologie de Berge pour les graphes ; mais si les termes de "sommet" et "arête" sont classiques, le rédacteur n'aime guère "sommet pendant" qui a l'air grotesque, et préférerait "point d'arrêt", qui irait assez bien avec "point de ramification".

Autre question de terminologie : le Congrès avait adopté d'enthousiasme une belle terminologie symétrique : "poids", "poids radiciels", "copoids", "copoids radiciels". Le malheur est qu'avec la définition adoptée, un "copoids radiciel" \uline{n'est pas} nécessairement un "copoids", ce qui est grammaticalement intolérable. Le rédacteur voit trois solutions possibles :

1) Prendre une autre terminologie pour "poids radiciels" et "copoids radiciels" ; c'est la solution adoptée dans la rédaction, qui dit respectivement "vocables" et "covocables", (dactylo, attention, ce n'est \uline{pas} "convocables"!). Le rédacteur ne tient pas spécialement à ces termes ; on pourrait aussi dire "noeuds" et "conoeuds" bien que ce dernier mot ait l'air un peu ... ce que vous pensez.

2) On ne peut naturellement changer le sens reçu de "racines" et "poids" ; mais "copoids" est beaucoup plus récent, (Chevalley, je crois), et pas encore très répandu ; on pourrait donc intervertir les sens de "copoids" et "\unsure{copóids|copoids} radiciels", les poids étant donc associés aux copoids radiciels, (et non plus aux copoids), et les copoids radiciels étant engendrés par les coracines, (ce qui est assez satisfaisant).

3) N'utiliser aucune terminologie pour "poids radiciels" et "copoids radiciels" ; l'expérience, (tout au moins dans la présente rédaction), montre qu'on s'en sert en fait très peu et qu'on s'en passe fort bien.
% archive sheet number: 6
\scanpage{8}\mspage{\handwritten{4}}
\hfill Bures le 12.8.19\deleted{\unsure{4|8|6}}64\\
Mon cher Serre,\\

Je commence à avoir de faibles lumières sur la façon de faire \handwritten{d}es VA avec les cycles al\handwritten{g}ébriquement équivalents à zéro d'une variété projective non singulière $X$ (sur un corps alg clos $k$). Si $X$ est connexe de dimension $n$, je sais associer à chaque entier $i$' \deleted{\unsure{ç|c}}compris entre $\deleted{\unsure{0|1}}1$ et $n$ une VA $J^i(X)$, jouant le rôle d'une "jacobienne intermédiaire" au sens de Weil, correspondant à des classes de cycles de codimension $i$ (pour une relation d'équivalence que je ne sais pas bien identifier pour l'instant, mais qui mérite certainement d'être appelée "équivalence d'Albanese"). Je suis moralement sûr que ce sont"les bonnes", bien que je n'ai pas prouvé grand-chose dessus pour l'instant; par contre, j'ai des conjectures en pagaille. Je te signale seulement que $J^1=\deleted{\unsure{B|P}}\operatorname{Pic}^{\unsure{o|0}}$ et $J^n=\operatorname{Alb}^{o}$, que $J^i$ et $J^{n+1-i}$ sont canoniquement duales l'une de l'autre, que $\dim J^i\handwritten{\leq\frac{1}{2}}b_{2i-1}$ (nombre de Bett\deleted{\unsure{a|e}}i) - de façon plus précise $T_{\unsure{\ell|l}}(J^i)$ est un quotient d'un sous-module de $H^{2i-1}(X,\underline{Z}_{\ell}(i))$ (et quand on sera plus savant, on saura prouver que c'est même un sous-module dudit), en caractéristique $0$ on peut également remplacer la cohomologie de Weil par la cohomologie de Hodge $H^i(X,\handwritten{\underline{\Omega}}_X^{i-1})$ ...

% slanted; two separate notes next to the Hodge formula
\margin{\handwritten{conjecturalement $\unsure{\underline{t}_{J^i}|t_{J^i}|\mathcal{Z}^{i}}\unsure{\subseteq|\subset} H^i(X,\underline{\Omega}_X^{i-1})$}}
\margin{\handwritten{\deleted{\unsure{$H^1(J^i)$|$H^1(J_i)$} $\simeq$ \unsure{Cl.|cl.} de Hodge}}}

Moyennant un certain nombre de choses non prouvées (qui risquent d'ailleurs d'être vaches) le théorème de positivité de Hodge se ramène à l'énoncé suivant, que je ne sais pas prouver, et que j'ai envie de te communiquer car il est bien terre à terre. Moralement, il s'agit de prouver que dans le cas où $\dim X=2m-1$, l'autodualité de $J^m$
% slanted
\margin{\handwritten{\unsure{canulé|annulé|Cancelé} tel quel, il faut tenir compte de la décomposition primitive\ldots\ldots}}
(qui s'exprime par une classe de correspondance divisorielle sur $J^m\mathbin{\mathrm{x}}J^m$ \deleted{anti}\uline{symétrique}, donc provenant - du moins modulo le facteur 2-
\scanpage{9}\mspage{\handwritten{5}}
d'un élément du groupe de Néron Severi de $J^m$) est \uline{positive} i.e. l'élément en question de Néron-Severi est ample, i.e. une polarisation. Pratiquement, cela s'explicite ainsi: Soit $T$ une variété de paramètres \inserted{connexe} non singulière munie d'un point marqué $a$ , $z$ une classe de cycles de codimension $m$ sur $T\mathbin{\mathrm{x}}X$ (à équivalence linéaire près, mettons), telle que $z(a)=0$ dans $X$, soi\deleted{\unsure{t}}ent $p$ et $q$ les deux projections de $T\mathbin{\mathrm{x}}T\mathbin{\mathrm{x}}X$ sur $T\mathbin{\mathrm{x}}X$, $r$ sa projection sur $T\mathbin{\mathrm{x}}T$, considérons
\[
D=r_{\handwritten{*}}(p^{\handwritten{*}}(z)\mathbin{.}q^{\handwritten{*}}(z))
\]
qui est une classe de diviseurs sur $T\mathbin{\mathrm{x}}T$, que nous considérons comme une classe de correspondance divisorielle sur $T\mathbin{\mathrm{x}}T$. Si $A=\operatorname{Alb}^{\unsure{o|0}}(T)$ (NB si tu veux, tu peux supposer $T=A$ et $a$ l'origine), elle provient donc d'une classe de correspondance sur $A\mathbin{\mathrm{x}}A$, évidemment symétrique. Soit $N$ le "noyau" de cette classe (i.e. le noyau de $A\handwritten{\longrightarrow}$ (duale de $A$) qu'elle définit), on obtient alors une classe de correspondance symétrique sur $J\mathbin{\mathrm{x}}J$ , où $J=A/N$ . A prouver que cette dernière est \uline{positive} ! Je me demande si les spécialistes "abéliens" pourraient avoir une idée sur une telle question, peut-être Matsusaka ? Ou toi-même ? Notes d'ailleurs que cette question te dévoile pratiquement la méthode de construction des $J^i$ généraux; \deleted{\unsure{s|S}}Si tu veux, tu peux te borner \inserted{aussi} au cas où tu disposes d'une sous-variété de codimension $m-1$ $Y$ de $X$, non singulière si tu y tiens, \deleted{\unsure{de|du}} et où $T=\operatorname{Pic}^{o}(Y)$, considéré comme paramétrant les diviseurs alg équiv à zéro de $Y$, mais considérés comme classes de cycles \deleted{dans} de codimension $m$ de $X$.

Merci pour la copie de la lettre à Ogg !

Bien à toi\\
% archive sheet number: 8
